\section{MPI+X}\label{sec:mpix}

\subsection{MPI+OpenMP: ranks versus threads}\label{sec:hybrid}
For a fixed number of cores, fewer ranks with more threads each means fewer and
larger shards: the ring has fewer steps and each step computes more, so communication
weighs less; in exchange, the threads of a rank must share its work evenly and a wide
rank spans more of the machine's memory hierarchy. \Cref{fig:hybrid} shows the
measured balance for $N=131{,}072$.

\begin{figure}[t]
\centering
\figureifexists{hybrid.pdf}
\caption{Run time for every split of (a) 72 cores of one node and (b) 288 cores of
four nodes into ranks $\times$ threads, $N=131{,}072$. Grey: the no-communication
control at the same split; bands show the range of three repetitions. Both axes
are logarithmic; the time axes cover only a few per cent, so small differences
look large.}
\label{fig:hybrid}
\end{figure}

Inside one node the split hardly matters: the seven splits of 72 cores lie within
\n{hybrid.range.pct.72}\% of each other, from \n{hybrid.best.time.72}\,s
(\n{hybrid.best.ranks.72} ranks $\times$ \n{hybrid.best.threads.72} thread) to
\n{hybrid.worst.time.72}\,s (one rank of 72 threads). Communication inside a node is
too cheap for fewer ranks to pay off, and the threaded kernel stays within that
margin of separate processes even as a single rank spanning both sockets: the
query tiles are independent, so threads share no state.
On four nodes fewer and larger ranks help a little. With 288 single-thread ranks
the run takes \n{hybrid.thin.time.288}\,s in this suite (the strong-scaling suite
of the same job measured \n{strong.T.big.288}\,s for the same configuration,
\cref{tab:strong}; the two differ by less than the protocol resolves); the best
split measured,
\n{hybrid.best.ranks.288} ranks of \n{hybrid.best.threads.288} threads, takes
\n{hybrid.best.time.288}\,s (\n{hybrid.gain.pct.288}\% less), a speedup of
\n{hybrid.speedup.288} over the serial kernel and the best result of this report
for $N=\n{hybrid.N.h}$. The direction agrees with \cref{sec:collapse}: with 32 ranks
a shard has \n{hybrid.c.best.288} rows instead of \n{strong.c.big.288}, and
$E_\text{comm}$ rises from \n{hybrid.ecomm.288.T1} to \n{hybrid.ecomm.288.T9}. The
gain is small because communication was a small part of the loss to begin with;
the clock, which no split changes, remains. It is also not monotonic: wider ranks
give it back (\n{hybrid.time.288.T36}\,s with 36 threads per rank,
\n{hybrid.time.288.T72}\,s with 72), and the control itself varies by a few per cent
between splits, since threads per rank also change how a rank's work meets the
sockets and caches. With repetitions that differ by up to
\n{hybrid.repspread.pct.288}\% for some splits, the ordering of neighbouring splits
is not resolved; what the data support is that a rank per quarter or half socket
is at least as good as 288 thin ranks on this machine for this problem.

\subsection{One GPU: a ladder of three kernels}\label{sec:gpu}
A GPU computes attention under the same block-at-a-time interface as the CPU core:
queries and the state $(m,\ell,a)$ stay on the device and one launch folds one
key/value block into the state. I wrote three kernels, each fixing the main weakness
of the one before; all are validated against the double-precision reference
(\cref{sec:validation}).

\begin{itemize}
\item \textbf{K0, naive.} One thread per (query, head) runs the serial two-pass row
loop. Every operand is read from global memory, the dot product is computed twice,
and each key updates the 64-element accumulator in global memory.
\item \textbf{K1, tiled.} Blocks of $32\times32$ threads compute $32\times32$ tiles
of scores from query and key tiles held in \term{shared memory} (the fast on-chip
memory of a thread block). The key tile is stored transposed, so that in the score
loop the 32 threads of a \term{warp} (the group of threads that executes together)
read 32 consecutive words, which lie in different memory banks and do not queue.
(Filling the transposed tile is not arranged that way: the threads of a warp then
write words that are 32 apart. I did not measure what that costs.) Scores are
written to global memory; a second kernel turns each row into weights; a third
multiplies weights and values tile by tile. The scores are thus \emph{materialised}:
one number per local query and key of the block, held in a scratch buffer of at
most 2\,GiB and processed in chunks of keys when a block does not fit.
\item \textbf{K2, fused.} One block owns 32 query rows of one head for the whole
key/value block. Each warp is one query row: its maximum and its sum are reductions
inside the warp, weights stay in shared memory, and each thread keeps two components
of the accumulator in registers until the block is finished. Nothing but the final
state is written to global memory. With a causal mask a tile of keys that lies
entirely in the future of the 32 queries is skipped before it is loaded.
\end{itemize}

\begin{figure}[t]
\centering
\figureifexists{gpu_ladder.pdf}
\caption{One NVIDIA A40, $d=d_v=64$, one head, no mask. (a) Run time and (b) useful
throughput of the three kernels; the tuned kernel on one CPU core is shown for scale.
``Useful'' counts $2(d+d_v)$ operations per allowed pair for every kernel, although
K0 computes each dot product twice. A GPU run warms up with one untimed repetition.}
\label{fig:gpuladder}
\end{figure}

\Cref{fig:gpuladder} shows the ladder. K0 reaches \n{gpu.k0.min} to
\n{gpu.k0.max}\,GFLOP/s, about the speed of one tuned CPU core: a direct port wastes
the device, because every operand comes from global memory and no thread shares
work with its neighbours. K1 reaches \n{gpu.k1.min} to \n{gpu.k1.max}\,TFLOP/s,
\n{gpu.k1.over.k0} times K0 at $N=\n{gpu.N.common.h}$: tiling through shared memory is
the decisive step. K2 is a further \n{gpu.k2.over.k1.small} to
\n{gpu.k2.over.k1.big} times faster, the more so the longer the sequence, and
reaches \n{gpu.k2.max}\,TFLOP/s. It computes $N=\n{gpu.k2.Nmax.h}$, that is
$\n{gpu.k2.pairs}\times10^{10}$ pairs, in \n{gpu.k2.Tmax}\,s without ever storing a
score outside the chip. These are the rates of this single-precision
implementation under this way of counting. They are \n{gpu.peakpct}\% of the
37.4\,TFLOP/s that the vendor states as the card's single-precision
peak~\cite{a40}, where the CPU kernel reaches \n{serial.peakpct}\% of its core's
peak: the GPU kernels are the less tuned of the two. I did not profile them
(occupancy, registers, memory throughput), so the numbers say what was reached,
not what limits it or how far the device could go.

\paragraph{The GPU against the CPU.} At $N=131{,}072$ one A40 with K2 needs
\n{gpu.k2.T131k}\,s. One tuned CPU core needs \n{strong.TS.big}\,s, a factor of
\n{gpu.vs.core.131072}; the 72 cores of one node need \n{gpu.node.T}\,s with the
overlapped ring. For this shape and precision, and for these two hand-written
implementations, one GPU therefore does the work of about \n{gpu.equiv.cores} cores
of a fully loaded node. The comparison is only meaningful because the CPU baseline
was tuned first: against the basic kernel the same GPU would have looked
\n{serial.speedup.tuned} times better. It is also lopsided the other way, since the
GPU kernel is further from its hardware limit than the CPU kernel; a better GPU
kernel would raise the equivalent core count and, by the second prediction of
\cref{sec:model}, make the exchange between GPUs matter at larger shards than
measured in \cref{sec:gpuring}.

\paragraph{Masks and real tensors.} With a causal mask K2 needs
\n{gpu.causal.k2.N65536} of its unmasked time at $N=65{,}536$, because key tiles
that lie in the future of a block of queries are skipped before they are loaded; K1
needs \n{gpu.causal.k1.N65536}: it skips the dot products of masked pairs, but it
still loads every tile, writes the whole block of scores and passes over it twice
more. On the real tensors of \cref{sec:real} ($N=32{,}768$, 14 query heads,
causal) K2 takes \n{gpu.real.T}\,s per layer and agrees with the double-precision
reference (\cref{tab:v7}).

\subsection{MPI+CUDA: the ring between GPUs}\label{sec:gpuring}
With one rank per GPU the ring of \cref{fig:ring} runs unchanged: a rank launches
the kernel on the block it holds and exchanges the next block meanwhile. The block
is either exchanged between host buffers and then uploaded (\term{host-staged};
nothing is copied back from the device, because keys and values are read-only) or
handed to a \term{CUDA-aware} MPI library as a device pointer. Each mode runs with
a blocking schedule (wait for the kernel, then exchange) and with an overlapped
one (exchange while the kernel runs).

\paragraph{Correctness.} The ring was checked at three levels. First, with two
ranks sharing one GPU, the allocation that was available first:
\n{gpucheck.configs} configurations covering both exchange modes and
schedules, the causal, window and sink+window masks, both partitions, the three
kernels and grouped heads pass on 64 sampled rows, with a largest error of
$\n{gpucheck.maxerr}$. That run exercises the code path; its timings mean nothing,
since both ranks compute on the same device. Second and third, with one rank per
GPU: two A40 in one node, and one A40 in each of two nodes. Each of these two jobs
contributes \n{gpu2.one.reps} repetitions, every one validated on
\n{gpu2.one.valrows} sampled rows, with largest errors of $\n{gpu2.one.maxerr}$ and
$\n{gpu2.two.maxerr}$. In the two-node job the script recorded what each node
showed: each rank saw exactly the one device that the scheduler had allocated.

\begin{figure}[t]
\centering
\figureifexists{gpu_ring.pdf}
\caption{Two A40 GPUs, fused kernel K2, $d=d_v=64$, one head, no mask: efficiency
$T_1/(2T_2)$ against rows per GPU, where $T_1$ is the time of one GPU measured in
the same job. Markers: median of three repetitions; bars: their range.}
\label{fig:gpuring}
\end{figure}

\paragraph{Two GPUs in one node.} A timed repetition contains the reset of the
state, the upload of the first block, all exchanges and the download of the
result. With the overlapped schedule two GPUs reach a speedup of
\n{gpu2.one.S.host.overlap.N131072} at $N=131{,}072$ (\n{gpu2.one.T1.N131072}\,s on
one GPU, \n{gpu2.one.T.host.overlap.N131072}\,s on two) and efficiencies between
\n{gpu2.one.E.overlap.min} and \n{gpu2.one.E.overlap.max} for $N\ge16{,}384$
(\cref{fig:gpuring}a). Most of what is missing is not communication. The time
inside the kernels alone gives nearly the same ratio
(\n{gpu2.one.Ek.host.overlap.N16384} at $N=16{,}384$ and
\n{gpu2.one.Ek.host.overlap.N131072} at $N=131{,}072$): the throughput of K2 grows
with the size of the block it is given (\cref{fig:gpuladder}b), and each of the two
GPUs works on blocks of half the size. This is the GPU's counterpart of the
granularity term in $E_0$.

The exchange becomes visible where \cref{sec:model} says it should, at few rows
per GPU and without overlap. Carrying the fitted CPU constant over to the GPU's
compute rate (\n{gpu.vs.loadedcore} times a loaded core) gives
$c^*\approx\n{gpu.cstar.pred}$ rows, hence an efficiency of \n{gpu.pred.E.c2048} at
2{,}048 rows per GPU if nothing overlapped. The blocking schedules measure
\n{gpu2.one.E.aware.blocking.N4096} (CUDA-aware) and
\n{gpu2.one.E.host.blocking.N4096} (host-staged) there, and the overlapped
host-staged schedule keeps \n{gpu2.one.E.host.overlap.N4096}. So the order of
magnitude of the extrapolation is right for a schedule that does not overlap, and
overlap removes most of the loss. It can do so here for a reason that does not
apply to CPU ranks: the kernel runs on the device and leaves the host thread free
to drive MPI. The CUDA-aware overlapped schedule is less regular at the two
smallest sizes (\n{gpu2.one.E.aware.overlap.N4096} and
\n{gpu2.one.E.aware.overlap.N8192}); these runs last a few milliseconds after a
warm-up of one repetition, and I do not read more into them.

\paragraph{One GPU in each of two nodes.} For large blocks nothing changes: the
speedup is \n{gpu2.two.S.host.overlap.N131072} at $N=131{,}072$, and the overlapped
schedule reaches \n{gpu2.two.E.overlap.min} to \n{gpu2.two.E.overlap.max} for
$N\ge16{,}384$ (\cref{fig:gpuring}b). For small blocks the exchange between the
nodes takes over. At $N=4{,}096$ one exchange takes
\n{gpu2.two.mpi.ms.host.overlap.N4096}\,ms where the kernel needs
\n{gpu2.two.kernel.ms.host.overlap.N4096}\,ms, and two GPUs are
\n{gpu2.two.slow.min.N4096} to \n{gpu2.two.slow.max.N4096} times \emph{slower} than
one. The blocking host-staged schedule pays about the same
\n{gpu2.two.mpi.ms.host.blocking.N16384}\,ms per exchange from $N=4{,}096$ to
$N=32{,}768$, which looks like a fixed delay of the transport and not a cost per
byte; I did not find its cause. These short configurations are also the least
repeatable ones of the report (bars in the figure). Two caveats belong here.
``CUDA-aware'' means that the library accepts device pointers; which path the
bytes take between the nodes was not determined. And the job log shows that the
library's components for collective operations failed to connect between the two
nodes and fell back to the default ones; the point-to-point path that the ring
uses worked, and every repetition validated.

\paragraph{Causal mask on two GPUs.} With two ranks a contiguous partition gives
the second GPU three quarters of the pairs and the zigzag partition gives each
one half. At $N=65{,}536$ the contiguous run takes
\n{gpu2.one.causal.T.contiguous}\,s and the zigzag run
\n{gpu2.one.causal.T.zigzag}\,s in one node (\n{gpu2.two.causal.T.contiguous} and
\n{gpu2.two.causal.T.zigzag}\,s across two): \n{gpu2.one.causal.ratio.contiguous} and
\n{gpu2.one.causal.ratio.zigzag} of the unmasked time, close to the
$\tfrac34$ and $\tfrac12$ that the pair counts give.
